\chapter*{Chapitre 11 — Études de cas et dimensionnement}
\addcontentsline{toc}{chapter}{Chapitre 11 — Études de cas et dimensionnement}

% ============================================================
\section*{Énoncés}
% ============================================================

% ----------------------------------------------------------
\exercice{Dimensionnement chauffage maison individuelle RE2020}
% ----------------------------------------------------------

\textbf{Contexte :}
Une maison individuelle de plain-pied est construite à Strasbourg (zone H1c,
$T_{ext,base} = \SI{-11}{\celsius}$). La surface habitable est de \SI{120}{\metre\squared}
pour une hauteur sous plafond de \SI{2{,}5}{\metre}. La maîtrise d'ouvrage souhaite
une installation conforme RE2020 avec une PAC air/eau alimentant un plancher chauffant basse température.

\textbf{Données de l'enveloppe :}

\begin{center}
\begin{tabular}{llcc}
  \toprule
  Paroi & Description & Surface (\si{\metre\squared}) & $U$ (\si{\watt\per\metre\squared\per\kelvin}) \\
  \midrule
  Murs extérieurs    & ITI laine roche 16 cm   & 140 & 0,20 \\
  Toiture            & Combles perdus 40 cm    & 120 & 0,12 \\
  Plancher bas       & Dalle isolée 10 cm      & 120 & 0,30 \\
  Vitrages           & Triple vitrage          &  28 & 0,80 \\
  Portes extérieures & Porte isolée            &   4 & 1,20 \\
  \bottomrule
\end{tabular}
\end{center}

\textbf{Ventilation :} VMC double flux avec récupérateur à plaques, rendement $\eta_{recup} = 0{,}88$.
Taux de renouvellement : $0{,}4\,\mathrm{vol/h}$.

\textbf{Température intérieure de consigne :} $T_{int} = \SI{19}{\celsius}$.

\bigskip
\textbf{Questions :}

\begin{enumerate}
  \item Calculer le coefficient de déperdition par transmission $H_T$ (\si{\watt\per\kelvin}).
  \item Calculer le débit volumique d'air neuf $q_v$ (\si{\metre\cubed\per\hour}) et le coefficient de déperdition par ventilation $H_V$ (\si{\watt\per\kelvin}) en tenant compte du récupérateur.
  \item En déduire la puissance de chauffage de pointe $\Phi_{ch}$ (\si{\kilo\watt}).
  \item La PAC choisie a un COP de $3{,}5$ à $-11/35\,\si{\celsius}$. Calculer la puissance électrique absorbée et la puissance installée en prenant une marge de 15\,\%.
  \item Calculer la puissance surfacique utile du plancher chauffant (\si{\watt\per\metre\squared}) et commenter le résultat par rapport aux valeurs courantes (40 à \SI{60}{\watt\per\metre\squared} max en RE2020).
  \item Estimer la consommation annuelle de la PAC (en kWh électriques) en utilisant la formule :
        $W_{el} = \dfrac{\Phi_{ch} \times DJU \times 24}{SCOP \times \Delta T_{base}}$
        avec $DJU = 3\,100$ et $SCOP = 4{,}0$.
  \item Calculer l'indicateur $Cep$ (en \si{\kilo\watt\hour_{ep}\per\metre\squared\per\year}) avec un coefficient de conversion $f_{prim} = 2{,}3$ et vérifier la conformité RE2020 (seuil résidentiel : \SI{70}{\kilo\watt\hour_{ep}\per\metre\squared\per\year} pour le seul chauffage).
\end{enumerate}

% ----------------------------------------------------------
\exercice{Dimensionnement d'une CTA pour locaux tertiaires}
% ----------------------------------------------------------

\textbf{Contexte :}
Un plateau de bureaux paysagers de \SI{500}{\metre\squared} (hauteur \SI{3}{\metre},
50 personnes) doit être équipé d'une centrale de traitement d'air (CTA) double flux.
Localisation : Lyon (zone H2b). Conditions d'été : $T_{ext} = \SI{33}{\celsius}$,
humidité relative extérieure = 50\,\%.

\textbf{Données :}
\begin{itemize}
  \item Débit d'air hygiénique : $25\,\si{\litre\per\second\per\personne}$
  \item Apports internes : personnes \SI{70}{\watt}/pers, éclairage \SI{8}{\watt\per\metre\squared}, bureautique \SI{20}{\watt\per\metre\squared}
  \item Apports solaires : vitrages Sud \SI{100}{\metre\squared}, $g = 0{,}30$, $G_{solaire} = \SI{550}{\watt\per\metre\squared}$
  \item Température de soufflage : \SI{15}{\celsius}
  \item Perte de charge réseau aéraulique circuit le plus défavorable : \SI{550}{\pascal}
  \item Rendement ventilateur : $\eta_{vent} = 0{,}70$
\end{itemize}

\bigskip
\textbf{Questions :}

\begin{enumerate}
  \item Calculer le débit d'air neuf $q_{v,AN}$ (\si{\metre\cubed\per\hour}).
  \item Calculer les apports internes totaux (\si{\kilo\watt}).
  \item Calculer les apports solaires (\si{\kilo\watt}).
  \item Calculer la charge due au renouvellement d'air (\si{\kilo\watt}).
  \item En déduire la puissance totale de la batterie froide $Q_{batt}$ (\si{\kilo\watt}).
  \item Sélectionner la température d'eau glacée nécessaire sachant que la batterie doit abaisser l'air de \SI{27}{\celsius} à \SI{15}{\celsius}. Quelle température d'entrée eau glacée est compatible ?
  \item Calculer la puissance du ventilateur de soufflage (\si{\kilo\watt}).
  \item Le groupe froid alimentant la batterie a un EER de $3{,}2$. Calculer sa puissance électrique absorbée.
  \item Calculer le débit d'eau glacée (\si{\metre\cubed\per\hour}) nécessaire pour la batterie, avec $\Delta T_{EG} = \SI{6}{\celsius}$.
\end{enumerate}

% ----------------------------------------------------------
\exercice{Analyse de performance d'une installation existante et préconisation}
% ----------------------------------------------------------

\textbf{Contexte :}
Un technicien de maintenance constate des anomalies dans le fonctionnement d'une installation
de climatisation d'un restaurant (superficie \SI{200}{\metre\squared}, 80 couverts).
Il dispose des mesures suivantes effectuées en conditions de pointe estivale :

\begin{center}
\begin{tabular}{lcc}
  \toprule
  Paramètre mesuré & Valeur relevée & Valeur attendue \\
  \midrule
  $T_{int}$ (salle) & \SI{29}{\celsius} & \SI{24}{\celsius} \\
  $T_{ext}$ & \SI{35}{\celsius} & — \\
  $T_{soufflage}$ CTA & \SI{22}{\celsius} & \SI{16}{\celsius} \\
  $q_v$ soufflage & \SI{2\,800}{\metre\cubed\per\hour} & \SI{3\,600}{\metre\cubed\per\hour} \\
  $T_{eau\ glacée\ entrée}$ batterie & \SI{12}{\celsius} & \SI{6}{\celsius} \\
  $T_{eau\ glacée\ sortie}$ batterie & \SI{15}{\celsius} & \SI{12}{\celsius} \\
  Courant moteur ventilateur & \SI{8}{\ampere} & \SI{9}{\ampere} \\
  Pression HP groupe froid & \SI{28}{\bar} & \SI{22}{\bar} \\
  Pression BP groupe froid & \SI{4}{\bar} & \SI{5}{\bar} \\
  \bottomrule
\end{tabular}
\end{center}

\textbf{Questions :}

\begin{enumerate}
  \item À partir des mesures, calculer la puissance frigorifique effectivement fournie par la batterie et la comparer à la puissance nominale (utiliser $\rho_{eau} = \SI{1000}{\kilogram\per\metre\cubed}$, $c_{p,eau} = \SI{4186}{\joule\per\kilogram\per\kelvin}$). On retiendra un débit d'eau glacée de $\SI{3{,}5}{\metre\cubed\per\hour}$ dans les deux cas.
  \item Identifier au moins 3 anomalies probables à partir des écarts constatés entre valeurs relevées et valeurs attendues.
  \item Pour chacune des anomalies identifiées, proposer une action corrective précise.
  \item La haute pression anormalement élevée (28 bar au lieu de 22 bar) sur un groupe froid au R-410A correspond à une température de condensation d'environ \SI{65}{\celsius} (au lieu de \SI{50}{\celsius} normale). Quel est l'impact sur le COP ? Justifier qualitativement.
  \item Le propriétaire envisage de remplacer le groupe froid R-410A (GWP = 2\,088) par un groupe R-32 (GWP = 675). Calculer la réduction d'équivalent CO\textsubscript{2} si la charge est de \SI{5}{\kilogram}. Commenter par rapport à la réglementation F-gas 2024.
\end{enumerate}

\newpage
% ============================================================
\section*{Corrigés}
% ============================================================

% ----------------------------------------------------------
\subsection*{Corrigé — Exercice 1 : Dimensionnement chauffage maison individuelle RE2020}
% ----------------------------------------------------------

\textbf{1. Coefficient de déperdition par transmission $H_T$}

\begin{center}
\begin{tabular}{lccr}
  \toprule
  Paroi & $U$ & $A$ & $U \cdot A$ (\si{\watt\per\kelvin}) \\
  \midrule
  Murs ext.   & 0,20 & 140 & 28,0 \\
  Toiture     & 0,12 & 120 & 14,4 \\
  Plancher bas& 0,30 & 120 & 36,0 \\
  Vitrages    & 0,80 &  28 & 22,4 \\
  Portes      & 1,20 &   4 &  4,8 \\
  \midrule
  \textbf{Total} & & & \textbf{105,6} \\
  \bottomrule
\end{tabular}
\end{center}

$$H_T = \SI{105{,}6}{\watt\per\kelvin}$$

\textbf{2. Coefficient de déperdition par ventilation $H_V$}

Volume du bâtiment : $V = 120 \times 2{,}5 = \SI{300}{\metre\cubed}$

Débit d'air neuf : $q_v = 0{,}4 \times 300 = \SI{120}{\metre\cubed\per\hour} = \SI{0{,}0333}{\metre\cubed\per\second}$

Avec récupérateur ($\eta = 0{,}88$) :
$$H_V = 0{,}34 \times q_{v,m^3/h} \times (1 - \eta_{recup}) = 0{,}34 \times 120 \times (1 - 0{,}88) = \SI{4{,}9}{\watt\per\kelvin}$$

\textbf{3. Puissance de chauffage de pointe}

$\Delta T = T_{int} - T_{ext,base} = 19 - (-11) = \SI{30}{\kelvin}$

$$H_{total} = H_T + H_V = 105{,}6 + 4{,}9 = \SI{110{,}5}{\watt\per\kelvin}$$
$$\Phi_{ch} = 110{,}5 \times 30 = \SI{3\,315}{\watt} \approx \SI{3{,}3}{\kilo\watt}$$

\textbf{4. Puissance PAC et puissance électrique}

Puissance électrique PAC nominale : $P_{el} = \dfrac{\Phi_{ch}}{COP} = \dfrac{3\,315}{3{,}5} = \SI{947}{\watt}$

Puissance installée avec marge de 15\,\% : $P_{installée} = 3{,}3 \times 1{,}15 \approx \SI{3{,}8}{\kilo\watt}$

Sélection : \textbf{PAC 4 kW} (gamme commerciale standard).

\textbf{5. Puissance surfacique plancher chauffant}

$$q_{plancher} = \frac{\Phi_{ch}}{A_{totale}} = \frac{3\,315}{120} = \SI{27{,}6}{\watt\per\metre\squared}$$

Valeur très confortable (largement inférieure à 40 W/m²), ce qui confirme que l'enveloppe très performante (RE2020) est compatible avec un plancher chauffant basse température à \SI{30}{\celsius} maximum de départ.

\textbf{6. Consommation annuelle PAC}

$$W_{el} = \frac{\Phi_{ch} \times DJU \times 24}{SCOP \times \Delta T_{base}} = \frac{3{,}315 \times 3\,100 \times 24}{4{,}0 \times 30} = \SI{2\,052}{\kilo\watt\hour\per\year}$$

\textbf{7. Indicateur $Cep$ et conformité RE2020}

$$Cep_{chauffage} = \frac{W_{el} \times f_{prim}}{A_{ref}} = \frac{2\,052 \times 2{,}3}{120} = \SI{39{,}3}{\kilo\watt\hour_{ep}\per\metre\squared\per\year}$$

$39{,}3 < 70$ : \textbf{conforme RE2020}. L'installation présente même un excellent niveau de performance énergétique, proche des exigences BEPOS.

% ----------------------------------------------------------
\subsection*{Corrigé — Exercice 2 : Dimensionnement d'une CTA pour locaux tertiaires}
% ----------------------------------------------------------

\textbf{1. Débit d'air neuf}

$$q_{v,AN} = 50\,pers \times 25\,\si{\litre\per\second} = \SI{1\,250}{\litre\per\second} = \SI{4\,500}{\metre\cubed\per\hour}$$

\textbf{2. Apports internes}

\begin{align*}
  Q_{personnes} &= 50 \times 70 = \SI{3\,500}{\watt} \\
  Q_{eclairage} &= 500 \times 8 = \SI{4\,000}{\watt} \\
  Q_{bureautique}&= 500 \times 20 = \SI{10\,000}{\watt} \\
  Q_{int,total}  &= \SI{17\,500}{\watt} = \SI{17{,}5}{\kilo\watt}
\end{align*}

\textbf{3. Apports solaires}

$$Q_{solaire} = g \times A_{vitrage} \times G_{solaire} = 0{,}30 \times 100 \times 550 = \SI{16\,500}{\watt} = \SI{16{,}5}{\kilo\watt}$$

\textbf{4. Charge de renouvellement d'air}

$\Delta T_{vent} = T_{ext} - T_{int} = 33 - 24 = \SI{9}{\kelvin}$ (température intérieure souhaitée : \SI{24}{\celsius})

$$Q_{vent} = \frac{1{,}2 \times 1005 \times 4500}{3600} \times 9 = 1507 \times 9 = \SI{13\,563}{\watt} \approx \SI{13{,}6}{\kilo\watt}$$

\textbf{5. Puissance batterie froide}

$$Q_{batt} = Q_{int} + Q_{solaire} + Q_{vent} = 17{,}5 + 16{,}5 + 13{,}6 = \SI{47{,}6}{\kilo\watt}$$

Avec marge de 10\,\% : $Q_{batt,installée} = 47{,}6 \times 1{,}10 \approx \SI{52}{\kilo\watt}$

\textbf{6. Température eau glacée}

Pour refroidir l'air de \SI{27}{\celsius} à \SI{15}{\celsius} (enthalpie sensible), la température d'entrée eau glacée doit être inférieure à la température de rosée de l'air traité. Avec un pincement de $\SI{4}{\celsius}$, $T_{EG,entrée} \leq 15 - 4 = \SI{11}{\celsius}$. On retient \textbf{\SI{7}{\celsius}} (circuit type 6/12\,°C).

\textbf{7. Puissance du ventilateur de soufflage}

$$P_{vent} = \frac{q_v \times \Delta P_{tot}}{\eta_{vent}} = \frac{(4500/3600) \times 550}{0{,}70} = \frac{1{,}25 \times 550}{0{,}70} = \SI{982}{\watt} \approx \SI{1}{\kilo\watt}$$

\textbf{8. Puissance électrique groupe froid}

$$P_{GF} = \frac{Q_{batt}}{EER} = \frac{52}{3{,}2} = \SI{16{,}25}{\kilo\watt}$$

\textbf{9. Débit eau glacée}

$$\dot{V}_{EG} = \frac{Q_{batt}}{\rho \cdot c_p \cdot \Delta T_{EG}} = \frac{52\,000}{1000 \times 4186 \times 6 / 3600} = \frac{52\,000}{6\,977} = \SI{7{,}45}{\metre\cubed\per\hour}$$

% ----------------------------------------------------------
\subsection*{Corrigé — Exercice 3 : Analyse de performance d'une installation existante}
% ----------------------------------------------------------

\textbf{1. Puissance frigorifique batterie — comparaison}

\textit{Puissance effectivement fournie (valeurs relevées) :}
$$Q_{batt,réel} = \dot{m} \times c_p \times \Delta T_{EG} = \frac{3{,}5}{3600} \times 1000 \times 4186 \times (15-12) = \SI{12\,227}{\watt} \approx \SI{12{,}2}{\kilo\watt}$$

\textit{Puissance nominale attendue :}
$$Q_{batt,nominal} = \frac{3{,}5}{3600} \times 1000 \times 4186 \times (12-6) = \SI{24\,454}{\watt} \approx \SI{24{,}5}{\kilo\watt}$$

La puissance réelle est donc à \textbf{50\,\% de la puissance nominale} — défaillance avérée.

\textbf{2. Anomalies identifiées}

\begin{enumerate}[label=\alph*.]
  \item \textbf{Température eau glacée trop élevée} ($12\,°C$ au lieu de $6\,°C$) : le groupe froid ne produit pas l'eau à la température requise. Cause probable : manque de fluide frigorigène, encrassement condenseur ou problème compresseur.
  \item \textbf{Débit d'air soufflé insuffisant} ($2\,800$ au lieu de $3\,600\,\mathrm{m^3/h}$) : obstruction réseau (filtre colmaté), ou perte de vitesse du ventilateur (courroie détendue, variateur mal réglé).
  \item \textbf{Haute pression anormalement élevée} ($28$ au lieu de $22\,\mathrm{bar}$) : condenseur encrassé (poussières, calcaire sur version eau), température ambiante du local technique trop élevée, ou ventilateur condenseur défaillant.
  \item \textbf{Basse pression trop faible} ($4$ au lieu de $5\,\mathrm{bar}$) : sous-charge en fluide frigorigène (fuite), détendeur mal réglé ou détendeur colmaté.
\end{enumerate}

\textbf{3. Actions correctives}

\begin{center}
\begin{tabular}{ll}
  \toprule
  Anomalie & Action corrective \\
  \midrule
  Eau glacée à 12°C & Contrôle charge fluide frigorigène ; recherche de fuite \\
  Débit d'air réduit & Nettoyage / remplacement filtre CTA ; réglage variateur \\
  HP élevée & Nettoyage condenseur (air/eau) ; vérification ventilateur condenseur \\
  BP basse & Test d'étanchéité circuit ; appoint fluide si fuite réparée \\
  \bottomrule
\end{tabular}
\end{center}

\textbf{4. Impact de la haute pression sur le COP}

Le COP d'une machine frigorifique est lié au rapport entre la puissance frigorifique et la puissance compresseur. Lorsque la température de condensation augmente (ici de 50°C à 65°C), la pression de condensation augmente, ce qui impose :
\begin{itemize}
  \item Un taux de compression plus élevé $\rightarrow$ puissance compresseur augmente ;
  \item Une enthalpie frigorifique utile qui diminue (diagramme de Mollier).
\end{itemize}
Le COP diminue donc significativement : une hausse de 15°C de la température de condensation peut entraîner une chute du COP de l'ordre de 20 à 30\,\%. Pour un EER nominal de 3,0, on peut estimer un EER réel de 2,1 à 2,4.

\textbf{5. Réduction d'équivalent CO\textsubscript{2} par changement de fluide}

\textit{R-410A (GWP = 2\,088), charge = 5 kg :}
$$CO_{2eq,R410A} = 5 \times 2\,088 = \SI{10\,440}{\kilogram_{CO2eq}}$$

\textit{R-32 (GWP = 675), charge = 5 kg :}
$$CO_{2eq,R32} = 5 \times 675 = \SI{3\,375}{\kilogram_{CO2eq}}$$

\textit{Réduction :}
$$\Delta CO_{2eq} = 10\,440 - 3\,375 = \SI{7\,065}{\kilogram_{CO2eq}} \approx \SI{7{,}1}{\tonne_{CO2eq}}$$

Soit une réduction de \textbf{67,7\,\%} de l'équivalent CO\textsubscript{2} potentiel (en cas de fuite totale).

La charge de 5 kg $\times$ GWP 2\,088 = 10,44 tonnes CO\textsubscript{2}eq : le seuil F-gas 2024 pour les nouveaux équipements est de \textbf{GWP $\leq$ 750}. Le R-410A (GWP = 2\,088) est donc \textbf{interdit pour les équipements neufs} depuis le 1\textsuperscript{er} janvier 2025. Le R-32 (GWP = 675) est conforme. Attention : le R-32 est légèrement inflammable (classe A2L) et nécessite des mesures de sécurité adaptées à l'installation et à la maintenance.
